% RESEARCH_DRAFT_SHA256: 24ac872cb843b0b6bdc6ae2a1316ac0bbbe49730ff178e45abb139cb16926fe4
% Layout and source-reference wrappers expanded for the web renderer.
\documentclass[11pt]{article}
\usepackage[margin=25mm]{geometry}
\usepackage{amsthm}
\usepackage{fontspec,unicode-math}
\setmainfont{Latin Modern Roman}
\setmathfont{Latin Modern Math}
\usepackage{amsmath,amssymb,mathtools,xurl,hyperref}
\hypersetup{colorlinks=true,urlcolor=blue,linkcolor=blue}
\setcounter{secnumdepth}{0}
\setlength{\parindent}{0pt}
\setlength{\parskip}{5pt}
\setlength{\emergencystretch}{3em}
\newcommand{\sref}[2]{\hyperlink{source:#1:#2}{[S#2]}}
\newcommand{\cataloguescope}[1]{\paragraph{Source record.}#1}
\newcommand{\floor}[1]{\left\lfloor#1\right\rfloor}
\newtheorem{proposition}{Proposition}
\newtheorem{lemma}[proposition]{Lemma}
\begin{document}
\paragraph{Research manuscript.} OpenAI \texttt{gpt-6-astra}, reasoning \texttt{ultra}; 11 October 2026.
\paragraph{Preserved source body.} The complete catalogue mathematical text and sources follow with layout and source links expanded, including source qualifications and its historical preparation notes. New research begins at the section headed Research attempt.

% UnsolvedMath ID 30005979; source catalogue code OWR-14298584-012
\section{Efficient Computation of Exceptional Galois Representations}\label{30005979}
{\small UnsolvedMath ID 30005979 · Computer Science}
\subsection{Definitions and mathematical statement}

\paragraph{Shared notation.}Standard number systems are $\mathbb N,\mathbb Z,\mathbb Q,\mathbb R,\mathbb C$. Specialized notation and index ranges follow the cited source and the definitions below.


A Galois representation is a homomorphism from the specified field\textquotesingle s Galois group into a linear group. Its projective representation identifies matrices differing by a nonzero scalar. Computing one requires a finite encoding of the relevant extension and representation data.
\paragraph{Record scope and qualifications.}
The catalogue title is retained as metadata. Named certificates, contracts, bounds or frameworks in that title are not asserted as results of the original source extract.
The exceptional representation family, input/output encoding, precision, and baseline complexity are not identified.

Source-record limitation: The exceptional Galois representations, input encoding, output requirements, and proposed improved complexity are not identified in the supplied algorithm-design question.
\paragraph{Statement recorded in the supplied source.}
Develop an efficient general algorithm for computing exceptional Galois representations, including their projective representations, with substantially better complexity than exhaustive class-field-theoretic or targeted number-field searches. \href{https://doi.org/10.4171/owr/2024/40}{[S2]}
\paragraph{Source status.}
The supplied catalogue labels this source record partially solved; that metadata was checked against the supplied snapshot on 2026-10-11. The label is preserved for provenance and is not an independent verification of present open status or mathematical correctness. The catalogue literature assessment is dated 2026-08-22.
\subsection{Short English statement}
Compute the cited exceptional Galois representations more efficiently.
\subsection{Sources}
\begin{itemize}
\item[S1] Ulam AI and the UnsolvedMath Contributors, UnsolvedMath: A Curated Collection of Open Mathematics Problems, record 30005979 (OWR-14298584-012). Catalogue curation and metadata: CC BY 4.0. The source extract is retained; definitions, English summaries and qualifications are editorial additions.
\url{https://huggingface.co/datasets/ulamai/UnsolvedMath}.
\item[S2] Explicit Methods in Number Theory (2025), Oberwolfach Reports, DOI 10.4171/owr/2024/40, pp. 2309--2368
\url{https://doi.org/10.4171/owr/2024/40}.
\item[S3] Booker, Artin's conjecture, Turing's method and the Riemann hypothesis
\url{https://arxiv.org/abs/math/0507502}.
\item Source-specific Oberwolfach Reports licences, including CC BY-SA 4.0 where applicable, are retained; see THIRD\_PARTY\_NOTICES.md.
\end{itemize}
\paragraph{Preparation and review.}
The original source and UnsolvedMath contributors retain attribution for the source extract. Codex (GPT-6; exact serving revision not exposed) assisted with definitions, English summaries, scope notes and formatting. No new human mathematical authorship or independently verified solution is asserted. Permanent source ID 30005979 and code OWR-14298584-012 checked against the complete supplied registry and current repository. Distinct upstream records, including repeated questions and broad or damaged extracts, are retained for complete coverage of the source CS category. A repeated record is not a new independent conjecture. Complete mathematical-equivalence and current-literature review has not been performed. Preparation and final review follow contribution prompts C and D, with source gaps disclosed. Structural and rendering checks do not establish a new theorem or open status.
\label{end:30005979}

\section{Research attempt}
\subsection{Literature, source context and scope}
The full original catalogue body and source qualifications are preserved above. The relevant primary workshop passages were read in the cached publisher report. Historical claims and subsequent attributed results are distinguished from the mathematical work below. A restricted example, standard lemma or known theorem is not counted as new progress on the entire original target. Missing assumptions and repeated or rhetorical extracts are not silently turned into different solved problems.

The primary report supplies structure missing from the short catalogue question: two-dimensional odd Artin representations with projective image $A_4,S_4$ or $A_5$, a conductor bound, and a GRH assumption for the described computations. In the $A_5$ stage it discusses Hunter searches and a squared-discriminant condition. This attempt tests whether the latter filter can replace identification of the Galois group.
\subsection{An irreducible square-discriminant false positive for an $A_5$ shortcut}
\begin{proposition}
The polynomial
\[
 f(x)=x^5+x^4-4x^3-3x^2+3x+1
\]
is irreducible over $\mathbb Q$, has square discriminant, and has Galois group cyclic of order five. Consequently irreducibility plus a square discriminant does not suffice to identify the required projective $A_5$ group.
\end{proposition}
\begin{proof}
Let $\zeta=e^{2\pi i/11}$ and $\theta=\zeta+\zeta^{-1}$. Put $s_j=\zeta^j+\zeta^{-j}$. The recurrence $s_{j+1}=\theta s_j-s_{j-1}$, starting with $s_0=2,s_1=\theta$, gives
\[
 s_2=\theta^2-2,\ s_3=\theta^3-3\theta,\ s_4=\theta^4-4\theta^2+2,\ s_5=\theta^5-5\theta^3+5\theta.
\]
The geometric-series identity $1+\sum_{j=1}^5s_j=0$ gives $f(\theta)=0$.

The polynomial $\Phi_{11}(x)=1+x+\cdots+x^{10}$ is irreducible: $\Phi_{11}(x+1)=((x+1)^{11}-1)/x$ is monic, all nonleading coefficients are divisible by 11, and its constant coefficient 11 is not divisible by $11^2$, so Eisenstein's elementary irreducibility criterion applies. Hence $[\mathbb Q(\zeta):\mathbb Q]=10$. The element $\zeta$ satisfies $z^2-\theta z+1=0$ and is nonreal whereas $\mathbb Q(\theta)\subset\mathbb R$, so the relative degree is two. Thus $\theta$ has degree five and $f$ is its minimal polynomial.

Every conjugate $\zeta^a+\zeta^{-a}$ is one of the recurrence polynomials $s_a(\theta)$, so all five distinct conjugates lie in $\mathbb Q(\theta)$. This field is the splitting field. The substitutions $\zeta\mapsto\zeta^a$ for $a\in(\mathbb Z/11\mathbb Z)^\times$ induce automorphisms with kernel $a=\pm1$ on $\theta$, giving a group of order five. Every group of prime order is cyclic, since any nonidentity element generates a subgroup whose order divides five.

A nonidentity permutation in this transitive order-five action is a 5-cycle and therefore even. The product of root differences $\prod_{i<j}(\theta_i-\theta_j)$ changes by the permutation sign under any automorphism, so it is fixed by the whole Galois group and is rational. Its square is the discriminant, proving that the latter is a rational square. The group of order five is not $A_5$, whose order is 60.
\end{proof}
\section{Verification and remaining gaps}
Exact recurrence and Sylvester-resultant computations verify $f$ and its discriminant $14641=11^4$. The group identification is proved above rather than inferred from that number. This does not criticize the source algorithm, which includes more information than the discarded filter. It gives neither efficient conductor enumeration, projective lifting nor a GRH-free algorithm for the exceptional representations.

\section{References and source attribution}
Primary workshop report: \url{https://doi.org/10.4171/owr/2024/40}. Relevant passages read: OWR 2024/40, printed pages 2353--2355, Allombert and Page.

All retrievals were checked on 11 October 2026. Original question attribution remains in the preserved Sources section. External theorems are attributed and not claimed as new discoveries or fully verified merely by citation.


\paragraph{Integration review.} Codex (GPT-6; exact serving revision not exposed) reviewed the written arguments and expanded layout and source-reference wrappers. Repository metadata and the production web renderer were checked using Python 3.13 and JavaScript with MathJax 3.2.2. This model-assisted review is not independent human proof verification; no PDF compilation or formal proof checking is asserted.

\section{Completion Estimate}
\noindent\textbf{COMPLETION: 0\%}
No substantiated new progress toward the entire original target is claimed. The worked results and failed route are delimited above; a known result, illustrative model, or auxiliary obstruction is not counted as original-target completion.

\section{Final status and attribution}
\textbf{Attempt status: unresolved.} This describes this research attempt, not a new ruling on the source catalogue's status or on external literature claims.
The mathematics and this self-audit were generated by OpenAI \texttt{gpt-6-astra}, reasoning \texttt{ultra}, on 11 October 2026. Codex (GPT-6; exact serving revision not exposed) orchestrated the batch, read and reviewed the written mathematical arguments, independently reran the finite checks, and contributed formatting/repository review. No human mathematical authorship, independent human proof checking, or formal verification is claimed. The historical preparation paragraphs in the reproduced catalogue body concern that body only.

Contribution prompt A was applied to the complete definition and supplied source record, followed by the adversarial checks of prompt B. Local reproduction files are \texttt{batch026/checks.py}, \texttt{computations.json}, and the per-ID audit JSON. They were not public links at preparation time. Numerical checks are finite and do not prove unrestricted results. No LaTeX engine was found locally; compilation was not performed by this research agent.

\end{document}
